%-*-latex-*-
\input{myassignmentpreamble}
\input{ciss358}
\input{yliow}
\renewcommand\TITLE{Assignment 7}

\newcounter{qc}
\newcommand\nextq{
  \newpage
  \addtocounter{qc}{1}
  Q{\theqc}.
  }


\begin{document}
\topmatter

\textsc{Objectives.}

\begin{enumerate}
  \li To solve problems using the dynamic programming method.
\end{enumerate}

Feel free to discuss among yourselves.
Implementations must be in C\texttt{++}.

As always, if an execution is shown, it is to fix the input and output
behavior of your program.
The output need not be correct.

In your main.cpp, at the top wrote down the asymptotic runtime
and space complexity:
\begin{Verbatim}[frame=single,fontsize=\small]
// T(n) = O(?)
// SPACE(n) = O(?)

... your program ...
\end{Verbatim}

Create directory \verb!a07!.
Put the latex files in \verb!a07!.
For the program for \verb!a07/a07q02/!
(main.cpp and other C\texttt{++} source files or data
files), put files in \verb!a07q02!.
Same for a07q03, etc.

SPRING2024:
For now, the only mathematical writing is in q01.tex.
For the other questions, you only need to implement C\texttt{++} program.
You are encouraged to try to write as many proofs for DP problems as you
can -- the proof of optimal substructure. After doing half a dozen you'll see
that they are all more or less the same.

%==============================================================================
\nextq
Prove that the rod cutting problem has the optimal substructure property.
Here the statement of the problem:x

\textsc{Rod cutting problem 1}.
Let $P[1..k]$ (\lq\lq prices'') be an array of positive values $>0$.
($P[i]$ is the price of a rod of length $i$. With this notation, you can
extend the array $P$ to index $0$ by
$P[0] = 0$ since the price of a rod of length 0 has to be 0.)
Given an integer $n > 0$, find $\ell_0, \ell_1, ..., \ell_{p-1}$
(the lengths of every cut; here, there are $p$ pipes) such that
\[
n = \ell_0 + \ell_1 + \cdots + \ell_{p - 1}, \,\,\,\,\,
1 \leq \ell_i \leq k
\]
and
\[
\sum_{i=0}^{p-1} P[\ell_i]
\]
is maximal.


\SOLUTION
\input{q01.tex}

%==============================================================================
\nextq
%-*-latex-*-
Write a C\texttt{++} program that solves the word segmentation problems.

Your program must include a trie (that's in CISS350).
A file (filename \verb!words.txt!)
is provided for your dictionary.
When loading each word in \verb!words.txt! into your dictionary,
convert all uppercase to lowercase and then ignore any word with non \verb!a-z! characters.
For instance ignore \verb!1st! if it's in the file.

\textsc{Test 1}
\begin{Verbatim}[frame=single,commandchars=\\\{\}]
\userinput{words.txt}
\userinput{toadsarenotfrogs}
2
to ads are not frogs
toads are not frogs
\end{Verbatim}
The first input is the name of the file of words.
The second input is the string to be segmented.
This is followed by the number of segmentations
and then all the segmentation; print a newline after a segmentation.
If there are more than one segmentation, print all of them.

(I mentioned before that all books on algorithms only
print one solution for a DP problem.
The word segmentation problem is an example of a DP problem where
you really should compute all solutions.)


%==============================================================================
\nextq
%-*-latex-*-
Adobe is experimenting with a new image compression
algorithm.
A color picture is an
$\verb!m! \times \verb!n!$
array \verb!P[0..m-1][0..n-1]! of pixels,
where each pixel specifies a triple of
red, green, and blue intensities.
To compress this picture slightly,
one pixel from each of the \verb!m! rows is removed,
so that the whole picture becomes one pixel
narrower, i.e., the new array is going to be
$\verb!m! \times \verb!(n - 1)!$.
To avoid disturbing visual effects, pixels
removed in two adjacent rows are in the same
or adjacent columns.
The collection of pixels removed is called a pixel-wide-cut.

To decide which pixels should be removed,
for each pixel
\verb!P[i][j]!,
an image processing algorithm
has calculated a real--valued
measure \verb!u[i][j]!,
indicating how bad the image file will be
if \verb!P[i][j]! is removed.
The lower the \verb!u[i][j]!, the
more similar the pixel is to its neighbors
therefore the better it would be if \verb!P[i][j]! is removed.
Removing \verb!P[i][j]! if \verb!u[i][j]! is high is bad.
The sum total of all the \verb!u[i][j]! for the pixels
\verb!P[i][j]!s in a pixel-wide-cut, is called the ugliness of that cut.
The goal is to minimize the ugliness.

(a)
How many pixel-wide-cuts are there for an
$\verb!m! \times \verb!n!$ image file?

(b)
Give an algorithm to find pixel-wide-cuts with the
minimal ugliness.

(c)
Implement your algorithm in C\texttt{++}.

(d)
What is the runtime and space complexity
of your algorithm? 
Write this in the documentation of your main.cpp:
\begin{Verbatim}[frame=single,fontsize=\small]
// T(n) = O(?)
// SPACE(n) = O(?)

... your program ...
\end{Verbatim}

Here's a sample output:
\begin{Verbatim}[frame=single,commandchars=\\\{\}]
\userinput{3 4}
\userinput{1.1 2.2 3.3 4.4}
\userinput{2.2 3.3 4.4 2.2}
\userinput{3.3 4.4 2.2 2.2}
6.6
((0,0), (1,0), (2,0))
((0,0), (1,1), (2,2))
\end{Verbatim}
The first input is the number of rows,
the second input is the number columns,
this is followed by the values of \verb!d!.
The first output, 6.6, is the minimal ugliness.
This is followed by the pixel-wide-cuts achieving the minimal ugliness.
In the above, there are two possible cuts:
the first is to remove pixels (0,0), (1,0), (2,0)
and the second is to remove pixels (0,0), (1,1), (2, 0).



%==============================================================================
\nextq
%-*-latex-*-
In order to transform one source string of text \verb!x[1..m]!
to a target string \verb!y[1..n]!,
we can perform various transformation operations.
Our goal is, given \verb!x! and \verb!y!,
to produce a series of transformations that change \verb!x! to \verb!y!.
We use an array \verb!z!
-- assumed to be large enough to hold all the characters it will need -- to hold
the intermediate results.
Initially, \verb!z! is empty, and at termination, we should have
\verb!z[j] = y[j]! for \verb!j = 1,2,...,n!.
We maintain current indices \verb!i! into \verb!x! and \verb!j! into \verb!z!,
and the operations are allowed to alter \verb!z! and these indices.
Initially, \verb!i = j = 1!.
We are required to examine every character in \verb!x! during the transformation, which
means that at the end of the sequence of transformation operations, we must have
\verb!i = m + 1!.

We may choose from among six transformation operations:

\begin{itemize}
  
\item
  Copy a character from \verb!x! to \verb!z! by setting \verb!z[j] = x[i]! and then incrementing both \verb!i!
  and \verb!j! . This operation examines \verb!x[i]!.

\item
  Replace a character from \verb!x! by another character \verb!c!, by setting \verb!z[j] = c!,
  and then
  incrementing both \verb!i! and \verb!j!.
  This operation examines \verb!x[i]!.

  \item
    Delete a character from \verb!x! by incrementing \verb!i! but leaving \verb!j! alone.
    This operation
    examines \verb!x[i]!.

  \item
    Insert the character \verb!c! into \verb!z! by setting \verb!z[j] = c! and then incrementing \verb!j!, but
    leaving \verb!i! alone.
    This operation examines no characters of \verb!x!.

  \item
    Twiddle (i.e., exchange) the next two characters by copying them from \verb!x! to \verb!z! but
    in the opposite order; we do so by setting \verb!z[j] = x[i + 1]! and \verb!z[j + 1] = x[i]!
    and then setting \verb!i = i + 2! and \verb!j = j + 2!. This operation examines \verb!x[i]! 
    and \verb!x[i + 1]!.

  \item
    Kill the remainder of \verb!x! by setting \verb!i = m + 1!.
    This operation examines all characters
    in \verb!x! that have not yet been examined.
    This operation, if performed, must
    be the final operation.
  \end{itemize}

  As an example, one way to transform the source string algorithm to the target
  string \verb!altruistic! is to use the following sequence of operations, where the
  underlined characters are \verb!x[i]! and \verb!z[j]! after the operation:


  \begin{Verbatim}[frame=single,commandchars=\\\{\}]
Operation         x            z
initial strings   \underline{a}lgorithm    _
copy              a\underline{l}gorithm    a_
copy              al\underline{g}orithm    al_
replace by t      alg\underline{o}rithm    alt_
delete            algo\underline{r}ithm    alt_
copy              algor\underline{i}thm    altr_
insert u          algor\underline{i}thm    altru_
insert i          algor\underline{i}thm    altrui_
insert s          algor\underline{i}thm    altruis_
twiddle           algorit\underline{h}m    altruisti_
insert c          algorit\underline{h}m    altruistic_
kill              algorithm_   altruistic_
\end{Verbatim}

Note that there are several other sequences of
transformation operations that transform
\verb!algorithm! to \verb!altruistic!.

Each of the transformation operations has an associated cost.
The cost of an operation depends on the specific application,
but we assume that each operation's
cost is a constant that is known to us.
We also assume that the individual costs of
the copy and replace operations are less than the combined costs of the delete and
insert operations; otherwise, the copy and replace operations would not be used.
The cost of a given sequence of transformation operations is the sum of the costs
of the individual operations in the sequence. For the sequence above, the cost of
transforming \verb!algorithm! to \verb!altruistic! is
\begin{align*}
  &3 \cdot \text{cost(copy)} + \text{cost(replace)} + \text{cost(delete)} + 4\cdot \text{cost(insert)} \\
  &\hspace{1cm} + \text{cost(twiddle)} + \text{cost(kill)}
\end{align*}

\begin{tightlist}
\item[(a)]
Given two sequences \verb!x[1..m]! and \verb!y[1..n]!
and set of transformation-operation
costs, the \defterm{edit distance} from x to y is the cost of the least expensive operation
sequence that transforms \verb!x! to \verb!y!. Describe a dynamic-programming algorithm
that finds the edit distance from \verb!x[1..m]! to \verb!y[1..n]! and prints an optimal operation
sequence. Analyze the running time and space requirements of your
algorithm.
Write a C++ program implementing your algorithm.
Here is a sample execution:
\begin{Verbatim}[frame=single,commandchars=\\\{\}]
\underline{algorithm}
\underline{altruistic}
\underline{1 2 3 1 2 5}
19
initial strings   >algorithm    
copy              a>lgorithm    a
copy              al>gorithm    al
replace by t      alg>orithm    alt
delete            algo>rithm    alt
copy              algor>ithm    altr
insert u          algor>ithm    altru
insert i          algor>ithm    altrui
insert s          algor>ithm    altruis
twiddle           algorit>hm    altruisti
insert c          algorit>hm    altruistic
kill              algorithm>    altruistic
\end{Verbatim}
The first two inputs are the strings \verb!x! and \verb!z!.
The next 6 inputs (integers) are the cost of copy, replace, delete, insert, twiddle, kill.
The first output is the cost.
The rest describe the operations.
\end{tightlist}

The edit-distance problem generalizes the problem of aligning two DNA sequences.
There are several
methods for measuring the similarity of two DNA sequences by aligning them.
One such method to align two sequences \verb!x! and \verb!y! consists of inserting spaces at
arbitrary locations in the two sequences (including at either end) so that the resulting
sequences \verb!x'! and \verb!y'! have the same length but do not have a space in the same
position (i.e., for no position \verb!j! are both \verb!x'[j]! and \verb!y'[j]! a space.) Then we assign a
\lq\lq score" to each position. Position \verb!j! receives a score as follows:
\begin{tightlist}
\li $+1$ if \verb!x'[j]! $=$ \verb!y'[j]! and neither is a space,
\li $-1$ if \verb!x'[j]! $\neq$ \verb!y'[j]! and neither is a space,
\li $-2$ if either \verb!x'[j]! or \verb!y'[j]! is a space.
\end{tightlist}
The score for the alignment is the sum of the scores of the individual positions. For
example, given the sequences \verb!x! $=$ \verb!GATCGGCAT! and \verb!y! $=$ \verb!CAATGTGAATC!, one
alignment is
\newpage
\begin{Verbatim}
    G ATCG GCAT
    CAAT GTGAATC
    -*++*+*+-++*
\end{Verbatim}
A \verb!+! under a position indicates a score of $+1$ for that position, a \verb!-! indicates a score
of $-1$, and a \verb!*! indicates a score of $-2$, so that this alignment has a total score of
$6 \cdot 1 - 2 \cdot 1 - 4 \cdot 2 = -4$.

\begin{tightlist}
\item[(b)]
Explain how to cast the problem of finding an optimal alignment as an edit
distance problem using a subset of the transformation operations copy, replace,
delete, insert, twiddle, and kill.
\end{tightlist}


%==============================================================================
\nextq
%-*-latex-*-
NSA is designing a new steganographic system to hide information.
Here's how it works.

Suppose agent Bob wants to send a message \verb!mad! to Alice
and Eve the rogue agent is eavesdropping.
Bob appends the characters of the message, except the last,
to the message but in the reverse order to get \verb!madam!.
Note that this is a palindrome with odd length.
He then inserts a huge number of characters into the string to further
obfuscate the message.
Here's an example:
\[
  \texttt{hieve\underline{\textred{m}}yh\underline{\textred{a}}mburger\underline{\textred{d}}oesnott\underline{\textred{a}}stegoodwith\underline{\textred{m}}armalade}
\]
The problem is Bob might accidentally create many subsequences in
this obfuscated message.
So Bob and Alice has agreed that if there is more than one
possible subsequence, the longest one is the correct palindrome.
NSA has asked you to write a program  
to find the longest odd-length palindromic subsequence
of a given input string (hopefully there's only one).
Of course it's possible that there are more than one longest odd-length
palindrome subsequence.
In that case, your program will be helpful in warning Bob that his obfuscation is ambiguous.

(a) For a string of length $n$, how many subsequences are there altogether?
(Using brute force, you can then check each subsequence to see if the
length is odd and is it is a palindrome. Of course you know by now that brute force is a bad idea.)

(b) Define a suitable measure for optimization (a length function)
in order to solve this problem.
For instance if the input is \verb!x[0..n-1]!, maybe
you want the function $L(\verb!i!)$ to be the length of the longest palindromic
subsequence of \verb!x[0..i]!, or maybe
you want the length of the longest palindromic
subsequence of \verb!x[i..n-1]!
or maybe you want
the function is $L(\verb!i!, \verb!j!)$ to be the length of the longest
palindomic subsequence of \verb!x[i..j]!.
Or maybe you want something that is completely different (and more
crazy than what I can think of.)

(c)
Find a recursion on your function in (b).
Implement a top-down DP solution.

(d)
Find a bottom-up DP solution to our problem
and implement it in C\texttt{++}.
Here's a sample run:
\begin{Verbatim}[frame=single,commandchars=\\\{\}]
\userinput{whatacharacter}
("carac", (5, 7, 8, 9, 10))
\end{Verbatim}
There is only one input. (Bob wanted to say \verb!car!, so the palindrome is \verb!carac!.)
The output is the longest odd-length palindromic subsequence followed by the
indices of this subsequence.
If there are more than one longest palindromic subsequences,
print them on separate lines.
For instance in the above, there's another output:
\begin{Verbatim}[frame=single,commandchars=\\\{\}]
\userinput{whatacharacter}
("carac", (5, 7, 8, 9, 10))
("hatah", (1, 2, 3, 4, 6))
\end{Verbatim}
(Note that for this part, you should print the longest.
\verb!tcaract! is actually longer. The above test case
is only to fix the input-output format. Outputs shown above need not be correct.)

(e) What is the runtime and space complexity of your algorithm?
Write this in the documentation of your main.cpp:
\begin{Verbatim}[frame=single,fontsize=\small]
// T(n) = O(?)
// SPACE(n) = O(?)

... your program ...
\end{Verbatim}

(f) Finally using the above program
write another program that prints all odd-length palindromic subsequence
with length of at least 3.
Organized the output in ascending order by length.
For string with the same length, organize them by dictionary order.

\textsc{Note.}
Assume characters in string are in lowercase.


%==============================================================================
\nextq
%-*-latex-*-
Polygonia is a high-end avant-garde architecture company.
Some of their designs are based on convex polygons.
For instance here's a convex polygon

\begin{python}
from latextool_basic import *
p = Plot()
p0 = (0,0)
p1 = (3,-1)
p2 = (4,2)
p3 = (2,3)
p4 = (-1,1)
p += Line(points=[p0,p1,p2,p3,p4,p0])
print(p)
\end{python}

\lq\lq Convex" means if you join two points of the polygon with a line segment,
the line segment is inside the polygon.
For course a fancy mansion has to be divided into rooms: bedrooms, banquet rooms,
study rooms, libraries, swimming pool rooms, cinema rooms, etc.
Polygonia is starting a new line of designs where
the rooms are created by building interior walls (walls inside the building)
joining vertices of a polygonal
building.
Here's an example where there are two interior walls for the above building:

\begin{python}
from latextool_basic import *
p = Plot()
p0 = (0,0)
p1 = (3,-1)
p2 = (4,2)
p3 = (2,3)
p4 = (-1,1)
p += Line(points=[p0,p1,p2,p3,p4,p0])
p += Line(points=[p0,p3], linecolor='red')
p += Line(points=[p0,p2], linecolor='red')
print(p)
\end{python}

and here's another:

\begin{python}
from latextool_basic import *
p = Plot()
p0 = (0,0)
p1 = (3,-1)
p2 = (4,2)
p3 = (2,3)
p4 = (-1,1)
p += Line(points=[p0,p1,p2,p3,p4,p0])
p += Line(points=[p1,p4], linecolor='red')
p += Line(points=[p1,p3], linecolor='red')
print(p)
\end{python}

Here's another wall design for a more complicated building:

\begin{python}
from latextool_basic import *
p = Plot()
p0 = (0,0)
p1 = (3,-1)
p2 = (4.5,0.8)
p3 = (4,2)
p4 = (2,3)
p5 = (0.5,2.5)
p6 = (-1,1)
p += Line(points=[p0,p1,p2,p3,p4,p5,p6,p0])
p += Line(points=[p1,p4], linecolor='red')
p += Line(points=[p1,p3], linecolor='red')
p += Line(points=[p4,p6], linecolor='red')
p += Line(points=[p4,p0], linecolor='red')
print(p)
\end{python}

In other words a design of interior walls is a collection of non--intersecting
line segments in the polygon building
such that each line segment joins two vertices of the polygon
and
the resulting rooms must be triangles.
The cost of interior walls is the sum of the lengths of the line segments
in the architectural design (think Pythagorus theorem).
The goal is to find the minimal cost for the interior
walls when given the vertices of the building.

(a) For a polygonal build with $n$ vertices,
how many different collections of interior wall design are there?
(SOLUTION PROVIDED.)

(b)
Find a suitable recursion that helps solve this problem.
(You have to find a function to optimize. Make sure you state the function clearly.)

(c)
Find and implement in C\texttt{++} a bottom-up DP solution to the problem.
Here's a sample execution
\begin{Verbatim}[frame=single,commandchars=\~\!\@]
~userinput!5@
~userinput!0.0 0.0 3.1 -1.0 4.2 2.0 2.1 3.1 -1.0 1.1@
12.7
{{0,2}, {0,3}}
\end{Verbatim}
The first input is the number of vertices of the polygon.
This is followed by the $(x,y)$--coordinates of the points listed in
counter-clockwise direction.
In the above, the polygonal is made up of points
$p_0 = (0.0, 0.0)$,
$p_1 = (3.1, -1.0)$,
$p_2 = (4.2, 2.0)$,
$p_3 = (2.1, 3.1)$, and
$p_4 = (-1.0, 1.1)$.
The output is 12.7 (the minimal cost of among all possible wall designs).
This is followed by the optimal wall design.
For instance in the above the first line wall is
\verb!{0,2}!, meaning the line segment $p_0, p_2$
and the second line segment is \verb!{0,3}! meaning $p_0, p_3$.
If there are two possible optimal wall designs, you must list all of them,
one complete wall design on each line.

(d) What is the runtime and the space complexity of your algorithm?
Write in the documentation of your main.cpp:
\begin{Verbatim}[frame=single,fontsize=\small]
// T(n) = O(?)
// SPACE(n) = O(?)

... your program ...
\end{Verbatim}

Solution and hint on the page next ...

\newpage
\textsc{WARNING ... incoming solution and hint}

(a)
Let $N(n)$ be the number of ways to triangulate an $n$--gon. Then
\[
N(n) = N(3)N(n-1) + N(4)N(n-2) + N(5)N(n-3) + \cdots + N(n-1)N(3)
\]
where $N(3) = 1$.
Let $N'(n) = N(n + 2)$.
Then the above becomes
\[
N'(n - 2) = N'(1)N'(n-3) + N'(2)N'(n-4) + N'(3)N'(n-5) + \cdots + N'(n-3)N'(1)
\]
Therefore
$N'(n - 2)$ is $C_{n-2}$, the $(n-2)$--th Catalan number.
Hence
\[
  N(n)
  =
  N'(n-2)
  = C_{n - 2}
  = \frac{1}{n - 2} \binom{2(n - 2 - 1)}{n - 2}
  = \frac{1}{n - 2} \binom{2(n - 3)}{n - 2}
\]
Asymptotically,
\[
  N(n)
  = C_{n - 2}
  = O(4^{n-2}/\sqrt{n-2})
\]


(c)
Consider this polygon as an example:
\begin{python}
from latextool_basic import *
p = Plot()
p0 = (0,0)
p1 = (3,-1)
p2 = (4.5,0.8)
p3 = (4,2)
p4 = (2,3)
p5 = (0.5,2.5)
p6 = (-1,1)

p += Line(points=[p0,p1,p2,p3,p4,p5,p6,p0])
#p += Line(points=[p1,p4], linecolor='red')
#p += Line(points=[p1,p3], linecolor='red')
#p += Line(points=[p4,p6], linecolor='red')
#p += Line(points=[p4,p0], linecolor='red')

from latexcircuit import *

X = POINT(x=p0[0], y=p0[1], label='$p_0$', anchor='east'); p += str(X)
X = POINT(x=p1[0], y=p1[1], label='$p_1$', anchor='north'); p += str(X)
X = POINT(x=p2[0], y=p2[1], label='$p_2$', anchor='west'); p += str(X)
X = POINT(x=p3[0], y=p3[1], label='$p_3$', anchor='south west'); p += str(X)
X = POINT(x=p4[0], y=p4[1], label='$p_4$', anchor='south'); p += str(X)
X = POINT(x=p5[0], y=p5[1], label='$p_5$', anchor='south east'); p += str(X)
X = POINT(x=p6[0], y=p6[1], label='$p_6$', anchor='east'); p += str(X)

print(p)
\end{python}
Let $C(0, 6)$ to be the cost of the polygon.
If I divide the polygon this way:
\begin{python}
from latextool_basic import *
p = Plot()
p0 = (0,0)
p1 = (3,-1)
p2 = (4.5,0.8)
p3 = (4,2)
p4 = (2,3)
p5 = (0.5,2.5)
p6 = (-1,1)

p += Line(points=[p0,p1,p2,p3,p4,p5,p6,p0])
#p += Line(points=[p1,p4], linecolor='red')
p += Line(points=[p0,p3], linecolor='red')
p += Line(points=[p3,p6], linecolor='red')
#p += Line(points=[p4,p0], linecolor='red')

from latexcircuit import *

X = POINT(x=p0[0], y=p0[1], label='$p_0$', anchor='east'); p += str(X)
X = POINT(x=p1[0], y=p1[1], label='$p_1$', anchor='north'); p += str(X)
X = POINT(x=p2[0], y=p2[1], label='$p_2$', anchor='west'); p += str(X)
X = POINT(x=p3[0], y=p3[1], label='$p_3$', anchor='south west'); p += str(X)
X = POINT(x=p4[0], y=p4[1], label='$p_4$', anchor='south'); p += str(X)
X = POINT(x=p5[0], y=p5[1], label='$p_5$', anchor='south east'); p += str(X)
X = POINT(x=p6[0], y=p6[1], label='$p_6$', anchor='east'); p += str(X)

print(p)
\end{python}
I can consider the cost $C(0, 3)$ and $C(3, 6)$
where
$C(0, 3)$ is the cost for polygon $p_0 p_1 p_2 p_3$
and 
$C(3, 6)$ is the cost for polygon $p_3 p_4 p_5 p_6$.


Note that the cost $C(0,6)$ is the cost of the diagonals in the
$p_0...p_6$ polygon.
Likewise the cost $C(0,3)$ is the cost of the diagonals in
the $p_0 ... p_3$ polygon.
The cost of the line $p_0p_3$ is not included in $C(0,3)$.
Likewise the cost of the line $p_3 p_6$ is not in the cost $C(3,6)$.
Therefore cost of triangulating $p_0...p_6$ where $p_0p_3p_6$ is chosen to be in the
triangulation is
\[
  C(0,3) + C(3,6) + c(p_0,p_3) + c(p_3,p_6)
\]
where $c(p_i, p_j)$ is the cost of the wall from $p_i$ to $p_j$.
Of course we have to go through all the possible cases:
\[
  C(0, 6)
  =
  \min
  \{
  C(0,k) + C(k, 6) + c(p_0, p_k) + c(p_k, p_6)
  \mid
  k = 1, 2, 3, ..., 5
  \}
\]
In general
\[
  C(i,j)
  =
  \min
  \{
  C(i,k) + C(k, j) + c(p_i, p_k) + c(p_k, p_j)
  \mid
  i + 1 \leq k \leq j - 1
  \}
\]
It should be clear that
\[
  C(i, i + 1) = 0
\]

(e) The runtime is $O(n^3)$ and the space is $O(n^2)$.


\end{document}
